\subsection{window-\texorpdfstring{$6$}{6} 边界 sheet parity 的指数条件期望、crossed product 与冻结扇区读出}\label{subsec:conclusion-window6-sheet-parity-watatani-crossedproduct-freeze-capacity}
本小节在 window-\(6\) 的边界二层覆盖、\(\ZZ_2^2\)-中心扇区分解与预言机容量框架下，给出一组可直接调用的终态命题：其核心是将 sheet parity 的块内代数实现、冻结相高阶矩读出、相应的统计恢复率，以及跨分辨率门控阈值统一到同一可审计接口。

\begin{theorem}[window-\(6\) 边界 sheet parity 条件期望的有限 Watatani 指数]\label{thm:conclusion-window6-sheet-conditional-expectation-index}
令
\[
A_6:=\CC[\mathcal R_6^{\mathrm{bin}}]
\cong
\bigoplus_{x\in X_6}M_{d_6^{\mathrm{bin}}(x)}(\CC),
\]
并沿定理 \ref{thm:terminal-window6-geo-fixed-subalgebra-wedderburn} 与命题 \ref{prop:foldbin6-boundary-clifford-subalgebra} 的记号取边界集合 \(X_6^{\mathrm{bdry}}\subset X_6\)。
对每个 \(w\in X_6^{\mathrm{bdry}}\)，在有序基 \((V_6(w),V_6(w)+34)\) 下定义
\[
P_w:=\begin{pmatrix}1&0\\0&-1\end{pmatrix},
\qquad
X_w:=\begin{pmatrix}0&1\\1&0\end{pmatrix},
\]
并有 \(P_w^2=X_w^2=I\)、\(P_wX_w=-X_wP_w\)。

记
\[
p_{\mathrm{bdry}}:=\sum_{w\in X_6^{\mathrm{bdry}}}\widetilde P_w\in Z(A_6),
\]
其中 \(\widetilde P_w\) 为对应简单块中心幂等元。
对边界块定义
\[
B_w:=\{a\in M_2(\CC):P_waP_w=a\}\cong\CC\oplus\CC,
\]
并令
\[
B_6:=\Bigl(\bigoplus_{w\in X_6^{\mathrm{bdry}}}B_w\Bigr)\ \oplus\ (1-p_{\mathrm{bdry}})A_6.
\]
定义 \(E_{\mathrm{sheet}}:A_6\to B_6\) 为：在每个边界块上
\[
(E_{\mathrm{sheet}})_w(a):=\frac12(a+P_waP_w),
\]
在非边界块上为恒等。
则 \(E_{\mathrm{sheet}}\) 为忠实条件期望，且具有有限 Watatani 指数，并满足中心值指标元公式
\[
\Ind_{\mathrm W}(E_{\mathrm{sheet}})=1+p_{\mathrm{bdry}}
\quad\text{于 }Z(A_6).
\]
\leanverified{paper\_conclusion\_window6\_sheet\_conditional\_expectation\_index}
\end{theorem}
\begin{proof}
在边界块 \(M_2(\CC)\) 内按 \(P_w\)-分级写
\(
a=a_{\mathrm{ev}}+a_{\mathrm{odd}}
\)，其中
\(
P_wa_{\mathrm{ev}}P_w=a_{\mathrm{ev}}
\)、\(
P_wa_{\mathrm{odd}}P_w=-a_{\mathrm{odd}}
\)。
于是 \((E_{\mathrm{sheet}})_w(a)=a_{\mathrm{ev}}\)。

取
\[
u_{w,0}:=I,\qquad u_{w,1}:=X_w.
\]
由 \(X_wP_w=-P_wX_w\) 可知 \(X_w\) 交换奇偶分量，从而
\[
u_{w,0}(E_{\mathrm{sheet}})_w(u_{w,0}^\ast a)+u_{w,1}(E_{\mathrm{sheet}})_w(u_{w,1}^\ast a)
=a_{\mathrm{ev}}+X_w(E_{\mathrm{sheet}})_w(X_wa)
=a_{\mathrm{ev}}+a_{\mathrm{odd}}
=a.
\]
故 \(\{I,X_w\}\) 是边界块的 Watatani 准基，边界块指标元为
\[
u_{w,0}u_{w,0}^\ast+u_{w,1}u_{w,1}^\ast=I+X_wX_w^\ast=2I.
\]
非边界块上 \(E_{\mathrm{sheet}}=\mathrm{id}\)，准基可取 \(\{I\}\)，指标元为 \(I\)。
逐块拼接并在中心上求和，得到
\[
\Ind_{\mathrm W}(E_{\mathrm{sheet}})
=(1-p_{\mathrm{bdry}})\cdot 1+p_{\mathrm{bdry}}\cdot 2
=1+p_{\mathrm{bdry}}.
\]
条件期望与忠实性均由逐块定义直接给出。证毕。
\end{proof}

\begin{theorem}[边界块的 crossed product 刻画与 \texorpdfstring{\(\delta=34\)}{delta=34} 的几何实现]\label{thm:conclusion-window6-boundary-crossed-product-delta34}
\leanverified{paper\_conclusion\_window6\_boundary\_crossed\_product\_delta34}
对每个 \(w\in X_6^{\mathrm{bdry}}\)，令
\[
D_w:=B_w\cong \CC\oplus\CC,
\]
并定义交换两 sheet 的 \(\ast\)-自同构
\[
\theta_w(a):=X_waX_w.
\]
则存在规范 \(\ast\)-同构
\[
D_w\rtimes_{\theta_w}\ZZ_2\ \cong\ M_2(\CC),
\]
其将 crossed-product 的实现元送到 \(X_w\)，并把 \(D_w\) 识别为对角子代数 \(B_w\)。

此外，window-\(6\) 边界 uplift 满足
\[
(\Fold^{\mathrm{bin}}_6)^{-1}(w)=\{V_6(w),V_6(w)+34\},
\]
差值严格为 \(34\)，从而诱导无歧义 sheet parity（推论 \ref{cor:terminal-delta-parity-rule}）。
并且该翻转在几何稳定子 \(\sigma_{\mathrm{geo}}\) 的作用下实现（推论 \ref{cor:terminal-foldbin6-geo-stabilizer-z2}），故上式可视为 window-\(6\) 上“几何 \(\ZZ_2\) 交换边界两 sheet”的代数化表达。
\end{theorem}
\begin{proof}
由 \(X_w^2=I\) 且 \(X_w\) 共轭交换对角极小投影，\((D_w,X_w)\) 满足 crossed-product 的生成关系，故由泛性质得到满射
\[
\Phi_w:D_w\rtimes_{\theta_w}\ZZ_2\to M_2(\CC),
\]
且 \(\Phi_w|_{D_w}\) 为对角嵌入、\(\Phi_w(u)=X_w\)（\(u\) 为实现元）。
两端均为四维复 \(\ast\)-代数，故 \(\Phi_w\) 为同构。

规范性来自边界二层覆盖给出的有序基 \((V_6(w),V_6(w)+34)\)：在该基下 \(X_w\) 的矩阵形式正是“加 \(34\)”对应的 sheet 交换。
由推论 \ref{cor:terminal-window6-local-uplift-admissibility}，该非平凡位移在 window-\(6\) 中唯一几何可实现，故结论成立。证毕。
\end{proof}

\begin{theorem}[冻结相中高阶带符号 Watatani 矩对 \texorpdfstring{\(\chi_m\)}{chi\_m} 的指数收敛读取]\label{thm:conclusion-signed-watatani-moments-read-chi}
对每个分辨率 \(m\)，令
\[
A_{m,1}\cong\bigoplus_{x\in X_m}M_{d_m(x)}(\CC),
\]
并记中心对合与取向指纹
\[
\varepsilon_{\mathrm{scan}},\ \varepsilon_{\mathrm{rev}}\in Z(A_{m,1}),
\qquad
\chi(x)=\bigl(\operatorname{sgn}(c_x),\operatorname{sgn}(\iota_x)\bigr)\in\{\pm1\}^2
\]
如定理 \ref{thm:fold-groupoid-z2x2-central-idempotents}。
令 \(E_{m,1}:A_{m,1}\to Z(A_{m,1})\) 为分块归一化迹条件期望，\(\tau_{m,1}\) 为归一化迹，并记
\[
A_m^{\max}:=\{x\in X_m:d_m(x)=M_m\},
\qquad
M_m:=\max_{x\in X_m}d_m(x).
\]

对任意 \(q\in\ZZ_{\ge 0}\)、\(s,t\in\{0,1\}\)，有闭式
\[
\tau_{m,1}\!\left(\varepsilon_{\mathrm{scan}}^{\,s}\varepsilon_{\mathrm{rev}}^{\,t}\Ind_{\mathrm W}(E_{m,1})^{\,q}\right)
=\frac1{2^m}\sum_{x\in X_m}\operatorname{sgn}(c_x)^s\operatorname{sgn}(\iota_x)^t\,d_m(x)^{q+1}.
\]

进一步，设 \(a=q+1\) 且满足冻结相假设：存在 \(a_{\mathrm c}\) 与 \(\Delta(a)>0\)，使
\[
\pi_{m,a}(A_m^{\max})\ge 1-\exp\!\bigl(-m\Delta(a)+o(m)\bigr),
\qquad
\pi_{m,a}(x):=\frac{d_m(x)^a}{\sum_{y\in X_m}d_m(y)^a}
\]
（见定理 \ref{thm:fold-escort-groundstate-concentration}）。
并设最大纤维族在指纹上同质：存在 \(\chi_m=(\alpha,\beta)\in\{\pm1\}^2\) 使
\[
A_m^{\max}\subseteq\{x:\chi(x)=\chi_m\}.
\]
则归一化比值
\[
R_{s,t,q}(m):=
\frac{\tau_{m,1}\!\left(\varepsilon_{\mathrm{scan}}^{\,s}\varepsilon_{\mathrm{rev}}^{\,t}\Ind_{\mathrm W}(E_{m,1})^{\,q}\right)}
{\tau_{m,1}\!\left(\Ind_{\mathrm W}(E_{m,1})^{\,q}\right)}
\]
满足
\[
R_{s,t,q}(m)=\alpha^s\beta^t+O\!\bigl(\exp(-m\Delta(a)+o(m))\bigr),
\qquad m\to\infty.
\]
\end{theorem}
\begin{proof}
第一式即定理 \ref{thm:fold-groupoid-z2x2-joint-spectral-measure} 的带符号矩公式。
将 \((s,t)=(0,0)\) 的同式作为分母，可得
\[
R_{s,t,q}(m)
=\frac{\sum_x \operatorname{sgn}(c_x)^s\operatorname{sgn}(\iota_x)^t\,d_m(x)^a}
{\sum_x d_m(x)^a}
=\EE_{\pi_{m,a}}\!\left[\operatorname{sgn}(c_X)^s\operatorname{sgn}(\iota_X)^t\right].
\]
在 \(A_m^{\max}\) 上，被积函数恒等于 \(\alpha^s\beta^t\)。
于是
\[
\left|R_{s,t,q}(m)-\alpha^s\beta^t\right|
\le
2\,\pi_{m,a}(X_m\setminus A_m^{\max}).
\]
由冻结相集中估计
\(
\pi_{m,a}(X_m\setminus A_m^{\max})
\le \exp(-m\Delta(a)+o(m))
\)
即得结论。证毕。
\end{proof}

\begin{theorem}[\texorpdfstring{\(\chi_m\)}{chi\_m} 的变分刻画：中心扇区对高阶指数矩的唯一极大化]\label{thm:conclusion-chi-variational-sector-maximization}
\leanverified{paper\_conclusion\_chi\_variational\_sector\_maximization}
沿定理 \ref{thm:fold-groupoid-z2x2-central-idempotents}，记中心幂等元族
\(
\{e_{\alpha,\beta}\}_{(\alpha,\beta)\in\{\pm1\}^2}\subset Z(A_{m,1})
\)。
对 \(q\in\ZZ_{\ge0}\) 定义扇区权重
\[
W_{\alpha,\beta}(m,q):=
\tau_{m,1}\!\left(e_{\alpha,\beta}\Ind_{\mathrm W}(E_{m,1})^{\,q}\right)
=\frac1{2^m}\sum_{\chi(x)=(\alpha,\beta)}d_m(x)^{q+1}.
\]
在定理 \ref{thm:conclusion-signed-watatani-moments-read-chi} 的冻结相与最大纤维同质假设下，令 \(a=q+1>a_{\mathrm c}\)。
则当 \(m\) 充分大时，\((\alpha,\beta)=\chi_m\) 为 \(\{W_{\alpha,\beta}(m,q)\}\) 的唯一极大点，并且对任意 \((\alpha,\beta)\neq\chi_m\) 有
\[
\frac{W_{\alpha,\beta}(m,q)}{W_{\chi_m}(m,q)}
\le
\exp\!\bigl(-m\Delta(a)+o(m)\bigr).
\]
\end{theorem}
\begin{proof}
由 \(e_{\alpha,\beta}\) 的块投影定义，
\[
W_{\alpha,\beta}(m,q)
=\tau_{m,1}\!\left(\Ind_{\mathrm W}(E_{m,1})^{\,q}\right)\cdot
\pi_{m,a}\!\bigl(\{x:\chi(x)=(\alpha,\beta)\}\bigr).
\]
若 \((\alpha,\beta)\neq\chi_m\)，则该集合与 \(A_m^{\max}\) 不交，故
\[
\pi_{m,a}\!\bigl(\{x:\chi(x)=(\alpha,\beta)\}\bigr)
\le
1-\pi_{m,a}(A_m^{\max})
\le
\exp\!\bigl(-m\Delta(a)+o(m)\bigr).
\]
另一方面，
\(
\pi_{m,a}(\{x:\chi(x)=\chi_m\})\ge \pi_{m,a}(A_m^{\max})\to 1
\)。
代回即得所述比值界与唯一极大化。证毕。
\end{proof}

\begin{theorem}[冻结相中指纹扇区的严格译码律与总变差收缩]\label{thm:conclusion-window6-sheet-parity-strict-decoding-tv}
在定理 \ref{thm:conclusion-signed-watatani-moments-read-chi} 与定理 \ref{thm:conclusion-chi-variational-sector-maximization} 的设定下，固定 \(a=q+1>a_{\mathrm c}\)，并记
\[
\chi_m=(\alpha,\beta)\in\{\pm1\}^2.
\]
则对充分大的 \(m\) 有严格译码公式
\[
\chi_m=
\Bigl(
\operatorname{sgn}R_{1,0,q}(m),\,
\operatorname{sgn}R_{0,1,q}(m)
\Bigr).
\]

进一步，定义四扇区概率
\[
p_{\gamma,\delta}^{(m,q)}
:=
\frac{W_{\gamma,\delta}(m,q)}
{\sum_{(\gamma',\delta')\in\{\pm1\}^2}W_{\gamma',\delta'}(m,q)},
\qquad
(\gamma,\delta)\in\{\pm1\}^2,
\]
并记 \(p^{(m,q)}\) 为 \(\{\pm1\}^2\) 上由上述质量函数确定的概率测度。
则其到单点质量 \(\delta_{\chi_m}\) 的总变差距离满足
\[
\left\|p^{(m,q)}-\delta_{\chi_m}\right\|_{\mathrm{TV}}
\le
\frac{3\exp\!\bigl(-m\Delta(a)+o(m)\bigr)}
{1+3\exp\!\bigl(-m\Delta(a)+o(m)\bigr)}
\le
3\exp\!\bigl(-m\Delta(a)+o(m)\bigr).
\]
\end{theorem}
\begin{proof}
由定理 \ref{thm:conclusion-signed-watatani-moments-read-chi}，
\[
R_{1,0,q}(m)=\alpha+O\!\bigl(\exp(-m\Delta(a)+o(m))\bigr),
\qquad
R_{0,1,q}(m)=\beta+O\!\bigl(\exp(-m\Delta(a)+o(m))\bigr).
\]
由于 \(\alpha,\beta\in\{\pm1\}\)，当 \(m\) 充分大时，上述误差项绝对值均严格小于 \(1/2\)；
于是 \(R_{1,0,q}(m)\) 与 \(R_{0,1,q}(m)\) 的符号分别等于 \(\alpha\) 与 \(\beta\)，从而得到严格译码公式。
\par
令
\[
\eta_m:=\exp\!\bigl(-m\Delta(a)+o(m)\bigr).
\]
由定理 \ref{thm:conclusion-chi-variational-sector-maximization}，对任意 \((\gamma,\delta)\neq \chi_m\) 都有
\[
\frac{W_{\gamma,\delta}(m,q)}{W_{\chi_m}(m,q)}\le \eta_m.
\]
因此
\[
\sum_{(\gamma,\delta)\neq\chi_m}W_{\gamma,\delta}(m,q)
\le
3\eta_m\,W_{\chi_m}(m,q),
\]
从而
\[
p_{\chi_m}^{(m,q)}
=
\frac{W_{\chi_m}(m,q)}
{W_{\chi_m}(m,q)+\sum_{(\gamma,\delta)\neq\chi_m}W_{\gamma,\delta}(m,q)}
\ge
\frac{1}{1+3\eta_m}.
\]
于是
\[
1-p_{\chi_m}^{(m,q)}
\le
\frac{3\eta_m}{1+3\eta_m}.
\]
由于 \(\delta_{\chi_m}\) 是 \(\{\pm1\}^2\) 上的单点质量，故
\[
\left\|p^{(m,q)}-\delta_{\chi_m}\right\|_{\mathrm{TV}}
=
1-p_{\chi_m}^{(m,q)},
\]
代入上式即得结论。证毕。
\end{proof}

\begin{theorem}[冻结相中边界 sheet parity 的双矩恢复与指数样本律]\label{thm:conclusion-window6-sheet-parity-sampling-law}
\leanverified{paper\_conclusion\_window6\_sheet\_parity\_sampling\_law}
在定理 \ref{thm:conclusion-signed-watatani-moments-read-chi} 与定理 \ref{thm:conclusion-chi-variational-sector-maximization} 的设定下，固定 \(a=q+1>a_{\mathrm c}\)，并记
\[
\chi_m=(\alpha,\beta)\in\{\pm1\}^2
\]
为冻结相中唯一极大扇区的符号指纹。令
\[
U(x):=\operatorname{sgn}(c_x)\in\{\pm1\},
\qquad
V(x):=\operatorname{sgn}(\iota_x)\in\{\pm1\},
\qquad
x\sim \pi_{m,a}.
\]
取独立样本 \(X_1,\dots,X_N\sim \pi_{m,a}\)，并定义经验均值
\[
\overline U_N:=\frac1N\sum_{j=1}^{N}U(X_j),
\qquad
\overline V_N:=\frac1N\sum_{j=1}^{N}V(X_j).
\]
约定 \(\operatorname{sgn}(0)=1\)，并定义估计器
\[
\widehat\alpha_N:=\operatorname{sgn}(\overline U_N),
\qquad
\widehat\beta_N:=\operatorname{sgn}(\overline V_N).
\]
则存在 \(m_0(a)\) 与常数 \(C_a>0\)，使得对所有 \(m\ge m_0(a)\)，若定义
\[
\varepsilon_m:=C_a\exp\!\bigl(-m\Delta(a)+o(m)\bigr)<1,
\]
则有错误界
\[
\PP\!\bigl((\widehat\alpha_N,\widehat\beta_N)\neq(\alpha,\beta)\bigr)
\le
4\exp\!\left(-\frac{N}{2}(1-\varepsilon_m)^2\right).
\]
特别地，当 \(m\) 充分大以致 \(\varepsilon_m\le \frac12\) 时，
\[
\PP\!\bigl((\widehat\alpha_N,\widehat\beta_N)\neq(\alpha,\beta)\bigr)
\le 4e^{-N/8}.
\]
\end{theorem}
\begin{proof}
由定理 \ref{thm:conclusion-signed-watatani-moments-read-chi}，存在常数 \(C_a>0\) 使
\[
\bigl|\EE_{\pi_{m,a}}[U]-\alpha\bigr|
=
\bigl|R_{1,0,q}(m)-\alpha\bigr|
\le
\varepsilon_m,
\]
\[
\bigl|\EE_{\pi_{m,a}}[V]-\beta\bigr|
=
\bigl|R_{0,1,q}(m)-\beta\bigr|
\le
\varepsilon_m.
\]
于是
\[
\alpha\,\EE_{\pi_{m,a}}[U]\ge 1-\varepsilon_m,
\qquad
\beta\,\EE_{\pi_{m,a}}[V]\ge 1-\varepsilon_m.
\]

定义
\[
Y_j:=\alpha\,U(X_j)\in\{\pm1\},
\qquad
Z_j:=\beta\,V(X_j)\in\{\pm1\}.
\]
则 \(\EE[Y_j]\ge 1-\varepsilon_m\)、\(\EE[Z_j]\ge 1-\varepsilon_m\)。又由 \(\operatorname{sgn}(0)=1\) 的约定，事件 \(\widehat\alpha_N\neq \alpha\) 蕴含
\[
\overline Y_N:=\frac1N\sum_{j=1}^{N}Y_j\le 0,
\]
同理 \(\widehat\beta_N\neq \beta\) 蕴含
\[
\overline Z_N:=\frac1N\sum_{j=1}^{N}Z_j\le 0.
\]

由 Hoeffding 不等式，
\[
\PP(\widehat\alpha_N\neq \alpha)
\le
\PP(\overline Y_N\le 0)
\le
\exp\!\left(-\frac{N}{2}(1-\varepsilon_m)^2\right),
\]
\[
\PP(\widehat\beta_N\neq \beta)
\le
\PP(\overline Z_N\le 0)
\le
\exp\!\left(-\frac{N}{2}(1-\varepsilon_m)^2\right).
\]
对二者作并合界得到
\[
\PP\!\bigl((\widehat\alpha_N,\widehat\beta_N)\neq(\alpha,\beta)\bigr)
\le
2\exp\!\left(-\frac{N}{2}(1-\varepsilon_m)^2\right),
\]
从而特别有
\[
\PP\!\bigl((\widehat\alpha_N,\widehat\beta_N)\neq(\alpha,\beta)\bigr)
\le
4\exp\!\left(-\frac{N}{2}(1-\varepsilon_m)^2\right).
\]
当 \(\varepsilon_m\le \frac12\) 时，
\(
(1-\varepsilon_m)^2\ge \frac14
\)，
故
\[
\PP\!\bigl((\widehat\alpha_N,\widehat\beta_N)\neq(\alpha,\beta)\bigr)
\le
4e^{-N/8}.
\]
证毕。
\end{proof}

\begin{proposition}[边界 uplift 的最小外加寄存器阈值由预言机容量给出]\label{prop:conclusion-boundary-uplift-min-register-by-oracle-capacity}
\leanverified{paper\_conclusion\_boundary\_uplift\_min\_register\_by\_oracle\_capacity}
对任意分辨率 \(m\)，设
\[
C_m(B):=\sum_{x\in X_m}\min\bigl(d_m(x),2^B\bigr)
\]
为 \(B\)-比特预言机容量（定理 \ref{thm:fold-oracle-capacity-closed-form}）。
在边界子群胚上定义
\[
C_m^{\mathrm{bdry}}(B):=\sum_{w\in X_m^{\mathrm{bdry}}}\min\bigl(d_m^{\mathrm{bdry}}(w),2^B\bigr),
\qquad
d_m^{\mathrm{bdry}}(w):=\left|(\Fold_m^{\mathrm{bin}})^{-1}(w)\right|.
\]
则成立：
\begin{enumerate}
  \item 在 \(m=6\) 时，\(d_6^{\mathrm{bdry}}(w)=2\)（推论 \ref{cor:terminal-delta-parity-rule}），故
  \[
  C_6^{\mathrm{bdry}}(1)=\sum_{w\in X_6^{\mathrm{bdry}}}2
  =\left|\widetilde X_6^{\mathrm{bdry}}\right|,
  \]
  即 \(1\) 比特已足以实现边界满反演。
  \item 在 \(m\in\{7,8\}\) 时，边界纤维均为三元集合（表 \ref{tab:foldbin_boundary_lift_m6_m8}，亦见定理 \ref{thm:window6-bdry-sheet-parity-nonfunctorial-extension}），故
  \[
  C_m^{\mathrm{bdry}}(1)=\sum_{w\in X_m^{\mathrm{bdry}}}2
  =\left|\widetilde X_m^{\mathrm{bdry}}\right|-\left|X_m^{\mathrm{bdry}}\right|,
  \]
  每条边界纤维至少损失一个微态。
  因而任意 \(1\)-比特协议在边界上必非满反演，并推出满反演阈值 \(B_{\min}\ge 2\)。
\end{enumerate}
\end{proposition}
\begin{proof}
容量闭式按纤维逐项取 \(\min(d,2^B)\)，精确表示在 \(B\)-比特侧信息下每条纤维可区分的最大微态数。
当 \(d\le 2^B\) 时该纤维可全量反演；当 \(d>2^B\) 时至少损失 \(d-2^B\) 个微态。

对 \(m=6\)，边界纤维大小恒为 \(2\)，取 \(B=1\) 即逐纤维饱和。
对 \(m\in\{7,8\}\)，边界纤维大小恒为 \(3\)，取 \(B=1\) 时每纤维上限仅为 \(2\)，故总缺口恰为 \(|X_m^{\mathrm{bdry}}|\)。
因此 \(1\)-比特不可能满反演，而 \(2^B\ge 3\) 至少要求 \(B\ge 2\)。证毕。
\end{proof}

\begin{theorem}[window-\texorpdfstring{$6$}{6} 的经典闭包与量子膨胀由同一纤维场支配]\label{thm:conclusion-window6-classical-closure-quantum-dilation-same-fiber-law}
在 window-\(6\) 的已核验 bin-fold 口径下，设每一步微态更新后都引入最小无记忆闭包机制：条件于当前稳定类型
\[
X_t=w\in X_6,
\]
将微态立即均匀化到纤维
\[
F_w:=(\Fold_6^{\mathrm{bin}})^{-1}(w).
\]
记相应纤维多重度场为
\[
d_6^{\mathrm{bin}}(w):=\card{F_w}.
\]
则成立：
\begin{enumerate}
  \item 稳定类型过程闭合为 \(X_6\) 上唯一的可逆马尔可夫链，其转移核必为推送核 \(P_6\)；
  \item 其平衡测度必为
  \[
  \pi_6(w)=\frac{d_6^{\mathrm{bin}}(w)}{64};
  \]
  \item 由同一纤维多重度场 \(d_6^{\mathrm{bin}}\) 定义的折叠量子信道 \(\mathcal E_6\) 的最小 Kraus 秩，等价地最小 Stinespring 环境维数，为
  \[
  S^{\mathrm{bin}}_2(6):=\sum_{w\in X_6}\bigl(d_6^{\mathrm{bin}}(w)\bigr)^2.
  \]
\end{enumerate}
因此，同一个纤维多重度场 \(d_6^{\mathrm{bin}}\) 同时固定了 window-\(6\) 上唯一可接受的经典闭包与最小量子实现代价。
\end{theorem}
\begin{proof}
前两条正是结论 \ref{con:xi-window6-instant-thermalization-equivalence} 与定理 \ref{thm:terminal-foldbin6-pushforward-markov} 的内容：在强可并性失败的背景下，一阶无记忆闭合存在当且仅当引入“纤维内瞬时热化”，且一旦闭合存在，其转移核被唯一强制为
\[
P_6(w,w')=\frac{A_6(w,w')}{6\,d_6^{\mathrm{bin}}(w)},
\]
平衡测度被唯一强制为
\[
\pi_6(w)=\frac{d_6^{\mathrm{bin}}(w)}{64}.
\]

对第 \(3\) 条，定理 \ref{thm:op-algebra-fold-quantum-channel-minimal-environment-equals-s2} 表明：折叠量子信道的最小 Kraus 秩，等价地最小 Stinespring 环境维数，精确等于相应二阶碰撞矩。把 \(m=6\) 的 bin-fold 纤维场代入，便得到
\[
\min\dim(\text{Stinespring 环境})
=
S^{\mathrm{bin}}_2(6)
=
\sum_{w\in X_6}\bigl(d_6^{\mathrm{bin}}(w)\bigr)^2.
\]
故 classical closure 与 quantum dilation 的两端并非由两套独立数据控制，而是由同一纤维多重度场 \(d_6^{\mathrm{bin}}\) 同时支配。证毕。
\end{proof}

\begin{conjecture}[跨分辨率 sheet parity 延拓的最小门控位成本]\label{conj:conclusion-cross-resolution-sheet-parity-min-gating-bits}
记边界纤维多重度谱
\[
d^{\mathrm{bdry}}(m):=
\left\{
\left|(\Fold_m^{\mathrm{bin}})^{-1}(w)\right|
\,:\,
w\in X_m^{\mathrm{bdry}}
\right\}.
\]
定义“零外加门控口径下的 sheet parity 延拓”为：存在逐纤维无不动点的 \(\ZZ_2\)-作用，并与自然投影交换。

猜想对 dyadic bin-fold 家族有
\[
B_{\min}(m)=
\left\lceil\log_2\!\bigl(\max d^{\mathrm{bdry}}(m)\bigr)\right\rceil,
\]
并且
\[
B_{\min}(m)=1
\quad\Longleftrightarrow\quad
d^{\mathrm{bdry}}(m)\ \text{全为偶数}.
\]
其中：
\begin{enumerate}
  \item \(m=6\) 的可行性由二层 uplift 与 \(\delta=34\) 的规范性给出（推论 \ref{cor:terminal-delta-parity-rule}）；
  \item \(m\in\{7,8\}\) 的失败由边界三元纤维强制不动点现象给出（定理 \ref{thm:window6-bdry-sheet-parity-nonfunctorial-extension}），并与命题 \ref{prop:conclusion-boundary-uplift-min-register-by-oracle-capacity} 的容量缺口一致。
\end{enumerate}
\end{conjecture}
